\section{\graph 优化器}
\label{sec:optimizer}

对于每个经过验证的 \graph，\sys 的 {\em \graph 优化器}通过进一步执行{\em 布局优化}、{\em 运算符调度}和{\em 内存规划}来最大化其性能，如 \Cref{fig:overview} 所示。
\sys 将这些 \graph 优化推迟到验证之后，原因有以下两点。
第一，这些优化{\em 不}影响生成的 \graphs 的正确性；在生成 \graphs 时省略它们，可以缩小 \sys 需要考察的搜索空间，因为具有相同图拓扑但张量布局、运算符顺序或内存分配方案不同的 \graphs 在 \graph 生成器看来是相同的。
第二，在验证之后再应用这些优化，也缩小了这些优化的搜索空间，因为 \graph 优化器只需优化与输入在功能上等价的 \graphs。

\paragraph{张量布局。}\graph 优化器在核、块和线程层次上探索所有中间张量可能的数据布局，并选择最佳组合以最大化性能。我们将布局选择表述为一个约束优化问题，并使用整数线性规划（ILP）算法进行最优求解。
具体地，对于每个张量 $t$ 和 $t$ 的每种可能布局 $l$，引入一个布尔变量 $B_{t,l}$ 来表示张量 $t$ 是否使用布局 $l$。
核、块和线程层次的运算符可能对张量布局施加各种约束。
例如，要使用 cuBLAS 库~\cite{cublas} 的核进行矩阵乘法，两个输入张量的最内维度必须在最后两个维度中。这些限制被转换为关于 $B_{t, l}$ 的一系列线性约束。不同的张量布局可能导致不同的性能。例如，某些输入张量布局支持从设备内存到共享内存的批量拷贝，而其他布局不支持。\sys 引入一个代价函数来建模不同布局选择下每个运算符的性能。
\sys 使用现成的 ILP 求解器（即 Z3~\cite{z3}）来找到满足所有布局约束同时最小化代价的最优布局策略。

\paragraph{运算符调度。}在 \graph 中，有多种拓扑顺序可以执行运算符，不同的顺序可能产生不同的性能。
对于给定的输入 \graph，\graph 优化器通过最小化每个线程块内的线程级同步（即 CUDA 中的 {\tt \_\_syncthreads()}）来识别高效的运算符调度策略。
为实现这一目标，\sys 为每个节点标注一个{\em 深度}，定义为从任意输入运算符到该节点的最长路径长度。\sys 使用动态规划算法计算每个节点的深度，并按深度升序调度所有运算符。
这种方法最小化了生成的 CUDA 核所需的线程级同步次数，因为 \sys 只需在深度不同的运算符之间插入同步点。

\paragraph{内存规划。}第三类验证后优化是内存规划，用于确定核、块和线程层次所有中间张量的内存偏移量。\sys 将内存规划表述为一个{\em 动态存储分配}问题，并穷举所有可能的分配方案以发现最优策略。
